\documentclass[11pt]{article}

\usepackage[T1]{fontenc}
\usepackage[utf8]{inputenc}
\usepackage[margin=1in]{geometry}
\usepackage{setspace}
\usepackage[hidelinks]{hyperref}

\setlength{\parskip}{0.8em}
\setlength{\parindent}{0pt}

\begin{document}

\begin{center}
{\LARGE \textbf{Building African-Context Multimodal Research Data: Practical Lessons for Text and Data Mining}}\\[1em]
{\large Guy Gael Ishimwe Karekezi, Lynne Chepkwony, and Emile Lucky Muhigira}\\[0.5em]
Carnegie Mellon University Africa\\
\href{mailto:iguygael@andrew.cmu.edu}{iguygael@andrew.cmu.edu},
\href{mailto:lchepkwo@andrew.cmu.edu}{lchepkwo@andrew.cmu.edu},
\href{mailto:emuhigir@andrew.cmu.edu}{emuhigir@andrew.cmu.edu}
\end{center}

\section*{Abstract}
Text and data mining (TDM) is often associated with large text corpora, but many contemporary research problems require the analysis of mixed digital content, including image-text pairs shared on news and social media platforms. Our project emerged from a practical research need: to understand whether a benchmark-trained multimodal misinformation detector could be made more useful for African-context digital media. In doing so, we encountered issues that are directly relevant to universities and research-oriented institutions in Africa, especially around data access, dataset construction, annotation, infrastructure, and the realities of working with locally meaningful digital content.

The project began with Fakeddit, a benchmark dataset containing image-text examples and misinformation-oriented labels. This dataset was useful for building an initial lightweight research pipeline based on CLIP-derived multimodal features and logistic regression. However, once we applied the benchmark-trained model to African-context examples, benchmark performance no longer held. The model degraded noticeably on African-context data, especially on the misinformation cases most relevant to our research question. This motivated us to construct and curate a separate African-context dataset of image-text examples grounded in public scenes, civic content, sports, public order, health, and daily life.

The most important TDM-related lesson from this experience is that dataset availability is a major bottleneck. Benchmark datasets are valuable, but they rarely reflect the contextual diversity needed for African research questions. For our work, we had to rely on a combination of locally stored image-text examples, manual curation, and structured annotation to build a usable dataset. This process was time-intensive and showed how difficult it can be to obtain stable, reusable corpora that reflect local realities rather than dominant global online platforms.

A second lesson concerns annotation and research labor. Our dataset required careful human judgment to distinguish between image-text pairs that were semantically compatible and those that misleadingly reframed the image. The challenge was not only technical but interpretive, since some examples required contextual knowledge about public events, regions, or common misinformation patterns. This shows that TDM projects in Africa often depend on local knowledge, annotation capacity, and reproducible guidelines, not only on computational tools.

A third lesson concerns infrastructure and tooling. Even with a lightweight modeling pipeline, practical experimentation depended on stable storage, local image availability, notebook-based workflows, and enough compute to repeatedly extract multimodal features and evaluate models. For research institutions with constrained infrastructure, enabling TDM therefore requires more than access to algorithms; it also requires access to curated data, annotation workflows, storage organization, and manageable toolchains.

Our results showed that adding African-context training data improved the original benchmark-trained model on African examples while preserving strong benchmark performance. More importantly, the project demonstrates that African-context research often requires building or adapting data resources rather than relying only on imported benchmarks. We therefore see this work as a case study in the practical realities of TDM-oriented research in Africa, and as evidence of the need for stronger local datasets, annotation capacity, and sustainable research workflows.

\section*{Brief Biographies}
\textbf{Guy Gael Ishimwe Karekezi.} Guy Gael Ishimwe Karekezi is a Master's student in Engineering Artificial Intelligence at Carnegie Mellon University Africa. His work focuses on multimodal AI, misinformation detection, explainable AI, and practical machine learning systems for African contexts. He is particularly interested in adapting benchmark-trained models to underrepresented environments and building usable AI tools grounded in local realities.

\textbf{Lynne Chepkwony.} Lynne Chepkwony is a student at Carnegie Mellon University Africa with interests in applied artificial intelligence, data-centric machine learning, and socially grounded technology. Her work includes dataset curation, annotation, and research-oriented AI workflows, with a focus on building practical systems that respond to contextual challenges in African environments.

\textbf{Emile Lucky Muhigira.} Emile Lucky Muhigira is a student at Carnegie Mellon University Africa interested in machine learning, responsible AI, and practical data-driven systems. His work includes annotation, multimodal data preparation, and evaluation pipelines for applied AI research, with particular interest in how AI tools can be adapted to better support African use cases.

\end{document}
